\section{Limitations, claim status and availability}\label{sec:limits}

\paragraph{What the records are.}
Each record is an exact upper bound on the rank of $\MM nmp$ over $\F$, and nothing more: not a lower bound, not an optimality claim, and not a statement about any other ring. The schemes rely on characteristic two (the rank-47 scheme for $\MM444$ in particular) and are not known to lift to the integers or the rationals (no lift was attempted); they must not be advertised as records for those rings.
``Record'' means smaller than the entries of four tables at named revisions on 2026-09-29. It does not mean new to the literature. Under the claim policy of the project, in which a result that improves the project's baseline but has not been searched against the literature is labelled \emph{apparent}, that is the label we use, and the persistence check of \cref{sec:persistence}, which supersedes 87 of the 795 and matches 112 more, is a first step towards an audit, not the audit.

\paragraph{Most of the results are consequences of known constructions.}
Composition dominates (\cref{sec:nature}), and the constructions used are established: block sums, Kronecker products, block recombination with a rank-47 outer scheme \cite{sedoglavic2017seven,drevet2011optimization}, Perminov's serendipitous product \cite{perminov2026serendipitous}, coordinate projection. What is new is the application of these constructions to the schemes that are best in characteristic two, at a scale (the 5{,}456 formats up to $32\times32\times32$) that the tables did not cover over $\F$, and the exact record of each derivation. A lower rank for one leaf changes many records at once, so the list should be read as a snapshot of a compositional frontier, not as 795 independent discoveries; the prior-art check makes the point, since 199 of the 795 are superseded or matched by compositions and claims built from published schemes.

\paragraph{Comparators.}
The tables are moving, incomplete and of mixed field (\cref{sec:comparators}), and they omit GF(2)-valid prior art that is public (\cref{sec:persistence}). For 458 of the 733 strongest records no GF(2)-valid scheme is listed, so the bar is a rank whose field is unknown; conversely, where no GF(2)-valid scheme is listed, a scheme of ours above the listed rank is not a record even if it is the best known over $\F$, so the list undercounts. The scoring rules exclude commutative algorithms and never count a rational scheme as valid over $\F$; twenty of the 298 rational files have only odd denominators and would reduce modulo two, but none is below our rank for its format.
The records come from a single scoreboard at a single date; a re-scoring at a later date can lower or remove entries, as the July certificates show (23 of the 186 were no longer records by late September).

\paragraph{Coverage and search depth.}
We hold a verified scheme for 3{,}662 of the 5{,}456 formats (67.1\%): for every format with largest side at most 16, for 70\% of those with largest side 17--24 and for 59\% of those with 25--32. That a format is not a record says little about the tables' rank for it. The search was finite (about $3\times10^{11}$ walk proposals in the campaign, not counting the earlier stages and the fleet; \cite{peterson2026methods}); we make no claim that any rank is a local or global minimum of the search space, and we did not attempt formats with a side above 32.

\paragraph{Verification.}
Exactness is established directly, twice, by code from one author with a shared bit layout. The provenance is re-derived only in part (\cref{sec:verification}): 42 of 244 walks were replayed, five seeds and eight records have incomplete leaf provenance, and the steps of the committed certificates are documented but were not replayed.

\paragraph{Attribution and licensing.}
Every derived tensor keeps the credit of the schemes it starts from, recorded per row: Perminov's FastMatrixMultiplication (MIT), Witteveen's fmm-schemes (CC BY 4.0), the matmulcatalog of Lacelle, which carries no licence file, and the Lille catalogue of Sedoglavic, used for comparison only.
The project owner treats decompositions as mathematical objects that are redistributed with attribution; \todo{confirm with the matmulcatalog authors before publication, as their repository has no licence}.

\paragraph{Availability.}
The search and verification tools are in the package \texttt{tungsten-metaflip} (SPDX: MIT OR Apache-2.0 WITH LLVM-exception) of the Tungsten repository \cite{tungsten2026lang}, branch \texttt{metaflip-\allowbreak public-\allowbreak ternary-\allowbreak gf2} (the tools of the release are commit \texttt{00ffa04f9}; the sources of these papers, their scripts and the verification logs are in later commits of the branch) \todo{push and merge the branch before submission}.
The lineage bundles (which hold 691 of the 795 record tensors with their derivations), the record index with the best listed ranks and margins, the fleet checkpoints, the verification logs and a snapshot of the tools are release \texttt{evidence-2026-09-29} of \texttt{tungsten-lang/metaflip-archives}, currently private \todo{deposit publicly, for example on Zenodo, and cite the DOI}; the other 102 records are committed certificates of the repository, and the logs of the later Ruby check and of the closure rebuild are in the repository as well (\texttt{papers/verification}). The scripts of the prior-art check, the transcribed literature rows with their quotes, the per-record flags and the 185 schemes we rebuilt are described in \texttt{papers/prior-art} of the repository \todo{deposit the rebuilt schemes and the fetched texts with the release}. Appendix~\ref{app:all} lists every record with its format, rank, bars, class and lineage family.
